\appendices
\section{Directional Curvature along Any Optimizer's Step}
\label{app:math}

Two full gradients of the same loss and the parameter update between them suffice to estimate curvature along the update direction. The derivation applies to any optimizer because it starts from the step actually taken. Let $u=\theta_t-\theta_{t-1}\ne 0$, and assume that the fixed loss $\mathcal L$ is twice continuously differentiable along this segment. For full gradients, integrating the derivative of $\nabla\mathcal L(\theta_{t-1}+au)$ along the segment gives
\begin{equation}
 g_t-g_{t-1}=\int_0^1 H(\theta_{t-1}+au)u\,da.
 \label{eq:app-gradient-integral}
\end{equation}
Project this change onto the step direction and divide by the distance traveled. The resulting scalar is exactly the directional curvature averaged along the step,
\begin{equation}
 \frac{u^\top(g_t-g_{t-1})}{\|u\|^2}
 =\int_0^1\frac{u^\top H(\theta_{t-1}+au)u}{\|u\|^2}\,da.
 \label{eq:step-path-average}
\end{equation}
Computing the left-hand side requires only the two endpoint gradients and the applied step. The approximation enters when we use this average to estimate curvature at the starting point. If the Hessian is $L_H$-Lipschitz in operator norm along the step, the directional curvature changes by at most $aL_H\|u\|$ at fraction $a$ of the step. Integration gives
\begin{equation}
 \left|\frac{u^\top(g_t-g_{t-1})}{\|u\|^2}
       -\frac{u^\top H_{t-1}u}{\|u\|^2}\right|
 \le \tfrac12 L_H\|u\|.
 \label{eq:app-curvature-bound}
\end{equation}

A step crossing a nonsmooth boundary can violate the assumptions behind these identities even when both endpoint gradients exist.

\emph{Connection to GD and Adam.} For plain GD with nonzero consecutive gradients, substituting $u=-\eta g_{t-1}$ in Eq.~\eqref{eq:step-path-average} recovers $\widehat\kappa_t$ in Eq.~\eqref{eq:kappa-estimate}. For Adam, the same curvature estimate uses $u=-\eta_{t-1}d_{t-1}$ from Eq.~\eqref{eq:adam-update}. The raw-gradient diagnostic in Section~\ref{sec:adam} uses a different projection. For nonzero consecutive gradients, the cosine definition gives
\begin{equation}
 1-\rho_t\frac{\|g_t\|}{\|g_{t-1}\|}
 =-\frac{g_{t-1}^{\top}(g_t-g_{t-1})}{\|g_{t-1}\|^2}.
 \label{eq:adam-projected-response}
\end{equation}
Thus the diagnostic measures how the preceding update changed the component of the gradient along $g_{t-1}$. Under GD it equals $\eta\widehat\kappa_t$. Under Adam the applied step can point in a different direction, so this diagnostic alone generally cannot recover curvature along that step.

\emph{Mini-batch gradients.} If $g_t=\nabla\mathcal L(\theta_t)+\xi_t$, the measured quotient in Eq.~\eqref{eq:step-path-average} includes the additional term $u^\top(\xi_t-\xi_{t-1})/\|u\|^2$. For plain SGD this term reduces to $-\varepsilon_t/\eta$. The bound in Eq.~\eqref{eq:app-curvature-bound} controls the error from variation in curvature. Sampling introduces a separate source of error.

\emph{Loss change over the step.} For full gradients, the inner product of the applied step with the mean of its two endpoint gradients approximates the loss change. Integrating by parts gives the exact error,
\begin{equation}
 \begin{aligned}
 &L_t-L_{t-1}-\tfrac12(g_{t-1}+g_t)^\top u\\
 &\qquad=\int_0^1(\tfrac12-a)\,u^\top H(\theta_{t-1}+au)u\,da.
 \end{aligned}
 \label{eq:app-loss-remainder}
\end{equation}
The weight $\tfrac12-a$ integrates to zero, so subtracting the midpoint Hessian inside the integral leaves it unchanged. Under the same Lipschitz assumption, its magnitude is at most $L_H\|u\|^3/12$. For a quadratic loss the error vanishes, making the loss-rise condition in Eq.~\eqref{eq:loss-rise-alignment} exact for full-batch GD. For plain SGD, substituting the full endpoint gradients $g_{t-1}-\xi_{t-1}$ and $g_t-\xi_t$ and the applied step $u=-\eta g_{t-1}$ gives the three terms in Eq.~\eqref{eq:loss-curvature} up to this same remainder.
